\documentclass[11pt]{article}
\usepackage{icmlko}
\begin{document}

\title{사전 등록이 스크립트를 이겼다 --- 약한 채택 하나와 그것을 되돌릴 조건}
\author{월드모델 실험실}
\date{노트 504 · 2026년 8월 3일}
\maketitle

\begin{abstract}
노트 503 이 남긴 유일한 생존 후보 --- 세계애니 \texttt{grp} --- 를 짝 씨앗
$20$ 으로 확정했다.

\begin{center}
\begin{tabular}{lrrl}
\hline
자 & 차 & 부호 & $95\%$ 구간 \\
\hline
판 & $+0.0011$ & $15/20$ & $[+0.0003, +0.0019]$ \\
날짜 통제 판 & $+0.0004$ & $12/20$ & $[-0.0005, +0.0013]$ \\
무리 안 판 & $+0.0004$ & $11/20$ & $[-0.0006, +0.0014]$ \\
KR 만화 & $+0.0084$ & $16/20$ & $[+0.0049, +0.0119]$ \\
비게임 앱 & $-0.0020$ & $7/20$ & $[-0.0046, +0.0007]$ \\
\hline
\end{tabular}
\end{center}

사전 등록은 ``\textbf{앱이 $0$ 을 물고 나머지 넷이 양수면 채택}''이었다.
앱 구간이 $[-0.0046, +0.0007]$ 로 $0$ 을 물고 넷의 평균이 다 양수다 ---
\textbf{통과}.

\textbf{그런데 스크립트 안에 박아 둔 고정 문턱(집 밖 $-0.0010$)은
반대했다.} 같은 실행에서 ``조건 미달''과 ``채택''이 나란히 찍혔다.
\textbf{사전 등록이 이긴다} --- 그것은 실행 전에 쓴 것이고 문턱은 다른
후보를 위해 쓴 것이다. 이 불일치를 숨기지 않고 적는다.
\end{abstract}

\begin{figure}[t]
\centering
\includegraphics[width=\textwidth]{figs/f1.pdf}
\caption{\textbf{가} --- 자 다섯. \textbf{나} --- 앞 채택과 견줌.
\textbf{다} --- 채택 다섯의 크기.}
\end{figure}

\section{약하다는 것을 정확히 적는다}

\begin{itemize}
\item $0$ 밖인 자는 \textbf{판과 KR 둘뿐}이다. 날짜 통제 판 $\cdot$
무리 안 판 $\cdot$ 앱 셋은 $0$ 을 문다.
\item 판 이득 $+0.0011$ 은 바로 앞 채택(애니 트렌드 $+0.0045$)의
\textbf{$1/4$} 이고, 이 실험실이 두 해 동안 채택한 다섯 중 \textbf{제일
작다}(계절 $+0.0041$ · 특기 $+0.0049$ · 깊이 $12$ $-0.0032$ · 애니
$+0.0045$).
\item 부호 $15/20$ 은 배포 문턱($\geq 9/12 = 0.75$)을 \emph{정확히}
맞춘다. 넘지 않는다.
\end{itemize}

\textbf{그래도 넣는다.} 사전 등록이 통과했고, 노트 446 거부권(판이 오르는데
날짜 통제가 내려감)에 안 걸리고, 노트 473 거부권(판이 비기는데 무리 안이
내려감)에도 안 걸린다 --- 판이 $0$ 밖이므로 ``비긴다''가 아니다.

\begin{quote}
\textbf{미리 적는 되돌림 조건} --- 다음 사이클에 \textbf{짝 씨앗 $40$}
으로 다시 재서 \textbf{판 구간이 $0$ 을 물면 되돌린다.} 되돌린 상태의
이름은 ``\texttt{DOMDROP} 에 세계애니가 없는 F18''.
\end{quote}

\section{배선}

\texttt{DOMDROP = \{"애니": ("trend\_",), "세계애니": ("grp",)\}}. 가드 ---
세계애니 \texttt{grp} 값을 난수로 바꿔도 예보 최대 차
$0.00\mathrm{e}{+}00$, 애니(대조, \texttt{grp} 안 뺌)는
$2.28\mathrm{e}{+}02$. 판 $0.4854$(씨앗 $6$ 앙상블).

\section{이 사이클이 남긴 것}

노트 $497 \sim 504$ 여덟 노트가 ``도메인마다 뺄 축''이라는 한 갈래를 팠다.

\begin{center}
\begin{tabular}{ll}
\hline
얻은 것 & 잃은 것 \\
\hline
채택 둘(애니 트렌드 · 세계애니 \texttt{grp}) & 고르개 --- 없다(노트 501) \\
조항 여섯 & 노트 $497$ 상한 $+0.0161$ --- 소진(노트 502) \\
가드 둘(\texttt{\_split\_census} · \texttt{\_coverage\_shift}) & 조항 둘을 다음 노트에서 죽임 \\
\hline
\end{tabular}
\end{center}

판은 $0.4812 \to \mathbf{0.4854}$ 로 $+0.0042$ 올랐다. 같은 기간 노트
$490$ 이 잰 \textbf{자료 손잡이 하나가 $-0.0283$} 이다. \textbf{여덟
노트의 모형 작업이 자료 하나의 $15\%$ 다.}

\section*{남은 것}

\begin{itemize}
\item \textbf{씨앗 $40$ 재확인} --- 되돌림 조건.
\item \textbf{덮음 $\geq 1\%$ 로 후보 목록 다시 세우기}(노트 503).
\item \textbf{그리고 자료 쪽으로} --- 노트 $490 \cdot 491$ 이 가리킨 곳:
검색 시계열 재수집(절단 고침 뒤), 시장팝업 태거, CN 태거.
\end{itemize}

\end{document}
